\documentclass[12pt]{article}
\usepackage[T1]{fontenc}
\usepackage[utf8]{inputenc}
\usepackage{lmodern}
\usepackage{textcomp}
\usepackage{enumitem}
\usepackage{geometry}
\geometry{margin=1in}
\begin{document}
\begin{center}
Answer Key - Biology - Graduate Level Assessment 2 \
\small (Answers are ordered exactly as in the question paper.)
\end{center}

\section*{Multiple Choice}
\begin{enumerate}[label=\arabic*.]
\item \textbf{Answer:} A
\\ \textit{Options:} A) BMR can be measured by the duration and quality of sleep; B) BMR can be estimated by the level of physical activity in a daily routine; C) BMR can be determined by measuring heart rate; D) BMR can be determined by the amount of sweat produced during physical activity; E) BMR can be estimated by the frequency of urination
\\ \textit{Note:} The duration and quality of sleep can influence metabolic processes and hormonal regulation, thereby affecting the basal metabolic rate, as sleep patterns can impact energy expenditure and metabolic function.
\item \textbf{Answer:} A
\\ \textit{Options:} A) Apical tissue; B) Sclerenchyma tissue; C) Ground tissue; D) Vascular tissue; E) Cork tissue
\\ \textit{Note:} Apical tissue, located at the tips of roots and shoots, is responsible for new growth in plants as it contains meristematic cells that facilitate primary growth through cell division and differentiation.
\item \textbf{Answer:} C
\\ \textit{Options:} A) conjugation results in no genetic diversity; B) conjugation requires more than two organisms; C) all individuals are identical genetically, and conjugation could not occur; D) repeated fission in paramecia results in sterile offspring that cannot conjugate; E) the group of paramecia lacks the cellular structures necessary for conjugation
\\ \textit{Note:} Conjugation in paramecia requires genetic diversity to facilitate the exchange of genetic material, and if all individuals have descended from a single individual through repeated fission, they would be genetically identical, making conjugation unnecessary and impossible.
\item \textbf{Answer:} A
\\ \textit{Options:} A) Temperature, precipitation, soil chemistry, and the types of decomposers present; B) Only soil structure and chemistry; C) Precipitation, temperature, and the presence of specific predator species; D) Only temperature and precipitation; E) Only the concentration of trace elements
\\ \textit{Note:} The correct answer is supported by the understanding that temperature, precipitation, soil chemistry, and the presence of decomposers are fundamental abiotic and biotic factors that influence the survival, reproduction, and overall characteristics of organisms within an ecosystem.
\item \textbf{Answer:} D
\\ \textit{Options:} A) Heterozygous traits; B) Mendelian traits; C) Multiple alleles traits; D) Polygenic traits; E) Pleiotropic traits
\\ \textit{Note:} Polygenic traits are characteristics that are influenced by multiple genes, resulting in a continuous range of phenotypic variations, such as height or skin color.
\end{enumerate}

\section*{True / False}
\begin{enumerate}[label=\arabic*.]
\setcounter{enumi}{5}
\item \textbf{Answer:} True
\\ \textit{Options:} A) True; B) False
\\ \textit{Note:} Weekend hospitalization leads to delayed provision of intensive procedures and elevated 1-year mortality for elderly AMI patients. The existence of measurable differences in treatments raises questions regarding the efficacy of a single input regulation (e.g., mandated nurse staffing ratios) in enhancing the quality of weekend care. My results suggest that targeted financial incentives might be a more cost-effective policy response than broad regulation aimed at improving quality.
\item \textbf{Answer:} False
\\ \textit{Options:} A) True; B) False
\\ \textit{Note:} GP ablation did not prove to be beneficial for postoperative stable NSR. A complete left atrial lesion set and biatrial ablation are advisable for improving rhythm outcomes. Randomized controlled trials are necessary to confirm our findings.
\item \textbf{Answer:} True
\\ \textit{Options:} A) True; B) False
\\ \textit{Note:} Interns order significantly more arterial blood gases per infant than junior and senior residents on-call in the neonatal ICU. Additional study is required to see if the experience of housestaff is associated with a broader array of neonatal outcomes, such as morbidity and mortality.
\item \textbf{Answer:} False
\\ \textit{Options:} A) True; B) False
\\ \textit{Note:} The prevalence of PAD is high in nursing home residents. AAI is not predictive for IHD mortality in this population. In very frail elderly traditional risk factors and PAD are less important predictors of death compared to poor functional status, nutritional factors and previous cardiovascular disease.
\item \textbf{Answer:} False
\\ \textit{Options:} A) True; B) False
\\ \textit{Note:} The site of access in our study does not appear to influence complications specifically neural injury or recurrence rates.
\end{enumerate}

\section*{Long Form Response}
\begin{enumerate}[label=\arabic*.]
\setcounter{enumi}{10}
\item \textbf{Answer:} The removal of the macronucleus from a paramecium has profound physiological consequences, ultimately leading to cell death. The macronucleus is crucial for regulating the expression of genes necessary for the cell's day-to-day functions, including metabolism, growth, and reproduction. Without it, the paramecium cannot maintain homeostasis or perform essential cellular processes, such as protein synthesis and response to environmental stimuli. Additionally, the macronucleus plays a role in coordinating cellular activities, and its absence disrupts the overall cellular organization and function. Consequently, the inability to sustain vital functions results in diminished cell viability and eventual mortality.
\\ \textit{Note:} correct because the macronucleus is essential for regulating gene expression related to vital cellular functions, and its removal disrupts processes like metabolism and homeostasis, ultimately resulting in cell death.
\item \textbf{Answer:} Heritability in the context of a Drosophila melanogaster population, specifically regarding the trait of abdominal bristles, can be quantified using the data collected on this phenotype. By calculating the heritability ($h^2$) for abdominal bristle number, we find a value of 0.89, indicating a high genetic contribution to the observed variation in this trait. This suggests that the majority of the phenotypic variance is attributable to genetic differences rather than environmental factors. Consequently, such a high heritability coefficient implies that selection for bristle number would likely be effective in enhancing this trait in subsequent generations. Overall, the data reinforces the significant role of genetics in determining the phenotypic expression of bristle number in Drosophila melanogaster.
\\ \textit{Note:} correct because it accurately calculates a high heritability value of 0.89 for abdominal bristle number in Drosophila melanogaster, indicating that genetic factors predominantly influence the observed phenotypic variation in this trait.
\end{enumerate}

\vfill
\noindent\textit{End of Answer Key}
\end{document}